% Part: first-order-logic
% Chapter: proof-systems
% Section: tableaux

\documentclass[../../../include/open-logic-section]{subfiles}

\begin{document}

\iftag{FOL}
      {\olfileid{fol}{prf}{tab}}
      {\olfileid{pl}{prf}{tab}}

\olsection{\usetoken{P}{tableau}}

تستعمل ال!!{tableau}s ال!!{signed formula}s، حيث تستعمل أنظمة
!!{derivation} كثيرة ترتيبات من ال!!{sentence}s.
و!!^a{signed formula} زوج: أوله علامة قيمة صدق، $\True$ أو $\False$، وثانيه
!!a{sentence}، هكذا:
\[
\sFmla{\True}{!A} \text{ أو } \sFmla{\False}{!A}.
\]
وأما !!^a{tableau} فشجرة متفرعة إلى أسفل تُرتَّب فيها
!!{signed formula}s. أولها جملة \emph{فروض}، وبعدها
!!{signed formula}s تُستخرج من ال!!{signed formula}s المتقدمة عليها بقاعدة استدلال.
وبكل قاعدة نضيف إلى آخر الفرع واحدة أو أكثر من ال!!{signed formula}s، أو نضع اثنتين من
!!{signed formula}s متجاورتين. وفي هذه الصورة الثانية ينشطر الفرع، وتكون الصيغتان الموقّعتان المضافتان رأسي الفرعين الجديدين.

إذا أُعملت القاعدة في !!{signed formula} مركبة، كانت الزيادة
!!{signed formula}s ذات !!{formula}s فرعية مباشرة. ولكل علامة قاعدة، فتأتي القواعد أزواجًا.
فقاعدة $\TRule{\True}{\land}$ تعمل في~$\sFmla{\True}{!A \land !B}$،
وتجيز إضافة كلتا الصيغتين الموقّعتين
$\sFmla{\True}{!A}$ و~$\sFmla{\True}{!B}$ إلى آخر أي فرع يشتمل على
$\sFmla{\True}{!A \land !B}$. وأما
$\TRule{\False}{\land}$ فتشطر الفرع بوضع
$\sFmla{\False}{!A}$ و$\sFmla{\False}{!B}$ متجاورتين.
ويكون !!^a{tableau} مغلقًا متى اشتمل كل فرع على صيغتين موقّعتين متطابقتين صيغةً مختلفتين علامةً:
$\sFmla{\True}{!A}$ و$\sFmla{\False}{!A}$.

تُعرّف علاقة~$\Proves$ القائمة على ال!!{tableau}s كما يأتي:
$\Gamma \Proves !A$ إذا وفقط إذا وجدت مجموعة منتهية
$\Gamma_0 = \{!B_1, \dots, !B_n\} \subseteq \Gamma$ وكان ثمة
!!{tableau} مغلق للفروض
\[
\{\sFmla{\False}{!A}, \sFmla{\True}{!B_1}, \dots, \sFmla{\True}{!B_n}\} 
\]
فمثلًا، هذا !!{tableau} مغلق يبين أن $\Proves
(!A \land !B) \lif !A$:
\begin{oltableau}
  [\sFmla{\False}{(\formula{A} \land \formula{B}) \lif \formula{A}}, just = \TAss
    [\sFmla{\True}{\formula{A} \land \formula{B}}, just = {\TRule{\False}{\lif}[1]}
      [\sFmla{\False}{\formula{A}}, just = {\TRule{\False}{\lif}[1]}
        [\sFmla{\True}{\formula{A}}, just = {\TRule{\True}{\land}[2]}
          [\sFmla{\True}{\formula{B}}, just = {\TRule{\True}{\land}[2]}, close]
        ]
      ]
    ]
  ]
\end{oltableau}

تكون المجموعة~$\Gamma$ غير متسقة في حساب ال!!{tableau} إذا وجد
!!{tableau} مغلق للفروض
\[
\{\sFmla{\True}{!B_1}, \dots, \sFmla{\True}{!B_n}\} 
\]
لبعض~$!B_i \in \Gamma$.

ظهرت ال!!^{tableau}s في خمسينيات القرن العشرين عند إيفرت بِث وياكو هنتيكا، كلٌّ مستقل عن الآخر، ثم بسطها ريموند سموليان وأشاعها.
وهي يسيرة الاستعمال؛ فإن إنشاء !!a{tableau} إجراء شديد الانتظام. وهذا الانتظام نفسه في ال!!{tableau}s يعين على تنفيذها بالحاسوب.
غير أن قراءة !!a{tableau} قد تعسر، وقد تخفى علاقته بالبرهان. والطريقة واسعة التطبيق، حتى إن لأنساق منطقية كثيرة أنظمة
!!{tableau}. ومن فوائد ال!!^{tableau}s أيضًا استخراج
!!{structure}s ترضي ما يُعطى من !!{sentence}s أو مجموعات منها.
فإن أمكن إرضاء المجموعة امتنع وجود !!{tableau} مغلق لها؛ أي بقي في كل
!!{tableau} فرع مفتوح. وإذا استُنفدت على فرع مفتوح جميع القواعد الممكنة، أمكن أن تُقرأ منه ال!!{structure} المُرضية.
وللطريقة صلة وثيقة بحساب التتابعيات: فإن ال!!{tableau} المغلق في أصله
!!{derivation} مختصر من ذلك الحساب، لكن اتجاه كتابته معكوس.

\end{document}
